\begin{table}[H]

\caption{\label{tab:extra-eu-import-cost-impulses-country-2022}Import exposure and imported-input cost impulses by euro-area country, 2022}
\centering
\begin{tabular}[t]{l*{4}{>{\raggedleft\arraybackslash}p{1.45cm}}}
\toprule
 & \multicolumn{2}{c}{Energy products} & \multicolumn{2}{c}{Food products}\\
 & \multicolumn{2}{c}{Raw price shock: 83.9\%} & \multicolumn{2}{c}{Raw price shock: 21.0\%}\\
\cmidrule(lr){2-3}\cmidrule(lr){4-5}
Country & $M/x$ & $b$ & $M/x$ & $b$\\
\midrule
Austria & 0.79 & 0.66 & 0.11 & 0.02\\
Belgium & 2.36 & 1.98 & 0.24 & 0.05\\
Croatia & 2.62 & 2.20 & 0.24 & 0.05\\
Cyprus & 0.62 & 0.52 & 0.14 & 0.03\\
Estonia & 0.87 & 0.73 & 0.29 & 0.06\\
Finland & 2.27 & 1.91 & 0.13 & 0.03\\
France & 1.52 & 1.27 & 0.12 & 0.02\\
Germany & 1.20 & 1.01 & 0.10 & 0.02\\
Greece & 6.51 & 5.46 & 0.22 & 0.05\\
Ireland & 0.78 & 0.65 & 0.30 & 0.06\\
Italy & 2.00 & 1.68 & 0.29 & 0.06\\
Latvia & 0.73 & 0.61 & 0.23 & 0.05\\
Lithuania & 6.07 & 5.09 & 0.23 & 0.05\\
Luxembourg & 0.07 & 0.06 & 0.05 & 0.01\\
Malta & 0.81 & 0.68 & 0.10 & 0.02\\
Netherlands & 2.60 & 2.18 & 0.62 & 0.13\\
Portugal & 2.47 & 2.07 & 0.42 & 0.09\\
Slovakia & 2.96 & 2.48 & 0.16 & 0.03\\
Slovenia & 0.96 & 0.81 & 0.16 & 0.03\\
Spain & 2.70 & 2.27 & 0.29 & 0.06\\
\textbf{Euro area} & \textbf{1.80} & \textbf{1.51} & \textbf{0.21} & \textbf{0.04}\\
\bottomrule
\end{tabular}
\par\vspace{0.35em}\footnotesize\noindent Notes. The raw shocks are the common changes in the 2022 annual-average extra-EU27 import unit-value indices relative to 2021. $M/x$ is the extra-EU imported-input share of gross output, and $b=h(M/x)$ is the corresponding cost impulse; both columns are expressed in percent. Energy covers B and C19; food covers A01 and C10--12. The euro-area entries are ratios of EA20 sums. Sources: Eurostat international trade statistics, FIGARO, authors' calculations.
\end{table}
